% !TeX root = ../main.tex
\lecture{1}{YYYY-MM-DD}{Probability basics}
\label{lec:probability}

\noindent\emph{Template example: replace this material with your lecture notes.}

\subsection{Random variables and indicators}

Let $A$ be an event in a probability space. Its indicator is a random variable
that records whether $A$ occurs.
\begin{definition}[Indicator]
\label{def:indicator}
\[
    \mathbf{1}_{\{A\}} =
    \begin{cases}
        1 & \text{if $A$ occurs},\\
        0 & \text{otherwise}.
    \end{cases}
\]
\end{definition}
Its expectation equals the probability of the event:
\begin{equation}
    \mathbb{E}\left[\mathbf{1}_{\{A\}}\right] = \Pr[A].
    \label{eq:indicator-expectation}
\end{equation}

\subsection{A first concentration bound}

\begin{theorem}[Markov's inequality]
\label{thm:markov}
Let $X$ be a nonnegative random variable with finite expectation, and let $a>0$
be a threshold. Then
\begin{equation}
    \Pr[X\geq a] \leq \frac{\mathbb{E}[X]}{a}.
    \label{eq:markov}
\end{equation}
\end{theorem}
\begin{proof}
Pointwise, $a\mathbf{1}_{\{X\geq a\}}\leq X$. Take expectations, apply
Equation~\eqref{eq:indicator-expectation}, and divide by $a>0$.
\end{proof}

\begin{example}
If $\mathbb{E}[X]=2$, Theorem~\ref{thm:markov} gives $\Pr[X\geq 10]\leq 1/5$.
\end{example}
\begin{exercise}
Construct a nonnegative random variable for which the bound in the preceding
example is attained.
\end{exercise}

\subsection{Pseudocode}

Let $x_1,\ldots,x_n$ be a nonempty stream, where $n$ is its length. The index $t$
counts the items processed, $J$ is a fresh uniform random index, and $s$ is the
stored sample. Algorithm~\ref{alg:reservoir} keeps one sample in constant space.

\begin{algorithm}[htbp]
\caption{Reservoir sampling of one item}
\label{alg:reservoir}
\begin{algorithmic}[1]
    \Require Nonempty stream $x_1,\ldots,x_n$
    \Ensure A uniformly sampled stream position's value
    \State $s\gets x_1$
    \For{$t=2,\ldots,n$}
        \State Draw $J$ uniformly from $\{1,\ldots,t\}$
        \If{$J=1$}
            \State $s\gets x_t$
        \EndIf
    \EndFor
    \State \Return $s$
\end{algorithmic}
\end{algorithm}

\begin{remark}
Each of the $n$ stream positions has probability $1/n$ of being sampled. If a
value occurs at $k$ positions, its probability of being returned is $k/n$.
\end{remark}
